\chapter{Inversión}

\marginal{dos objetos, dos historias}
El capítulo empieza separando dos cosas que el género suele mezclar. El
\emph{gasto de inversión} de la clasificación económica incluye la inversión
financiera; la \emph{inversión física} no. En 2025 la diferencia entre las dos
lecturas es el capítulo entero: la inversión física baja de 3.0 a \textbf{2.3\%
del PIB} y la financiera de 0.7 a 0.4.

\begin{table}[htbp]
\centering
\caption{Inversión por capítulo de objeto del gasto (mdp)}
\label{tab:C8_1}
\small
\begin{tabular}{lSSSSS}
\toprule
{Concepto} & {2024} & {2025} & {Dif. nominal} & {Dif. real} & {Var. real \%} \\
\midrule
{Capitulo 6000 Inversion publica (obra) (GF + entidades)} & \num{447575.3} & \num{290693.9} & \num{-156881.4} & \num{-176127.1} & \num{-37.7} \\
{Capitulo 5000 Bienes muebles e inmuebles (GF + entidades)} & \num{78494.0} & \num{46296.1} & \num{-32198.0} & \num{-35573.2} & \num{-43.5} \\
{Capitulo 7000 Inversiones financieras y otras provisiones (GF + entidades)} & \num{340623.7} & \num{351922.6} & \num{11298.9} & \num{-3348.0} & \num{-0.9} \\
\multicolumn{6}{l}{\itshape Capitulo 6000 por ramo (GF)} \\
{13 Marina} & \num{0.0} & \num{22058.5} & \num{22058.5} & \num{22058.5} & {---} \\
{09 Infraestructura, Comunicaciones y Transportes} & \num{35330.5} & \num{11833.7} & \num{-23496.9} & \num{-25016.1} & \num{-67.9} \\
{16 Medio Ambiente y Recursos Naturales} & \num{35343.4} & \num{6682.1} & \num{-28661.3} & \num{-30181.0} & \num{-81.9} \\
{01 Poder Legislativo} & \num{261.4} & \num{326.7} & \num{65.3} & \num{54.1} & \num{19.8} \\
{12 Salud} & \num{174.9} & \num{49.2} & \num{-125.7} & \num{-133.2} & \num{-73.0} \\
{03 Poder Judicial} & \num{2.9} & \num{28.6} & \num{25.7} & \num{25.6} & \num{845.9} \\
{47 Entidades no Sectorizadas} & \num{100.6} & \num{24.3} & \num{-76.3} & \num{-80.7} & \num{-76.8} \\
{32 Tribunal Federal de Justicia Administrativa} & \num{3.1} & \num{3.1} & \num{0.0} & \num{-0.1} & \num{-4.1} \\
{06 Hacienda y Crédito Público} & \num{25.1} & \num{0.0} & \num{-25.1} & \num{-26.2} & \num{-100.0} \\
{07 Defensa Nacional} & \num{92671.5} & \num{0.0} & \num{-92671.5} & \num{-96656.3} & \num{-100.0} \\
{11 Educación Pública} & \num{1170.0} & \num{0.0} & \num{-1170.0} & \num{-1220.3} & \num{-100.0} \\
{43 Instituto Federal de Telecomunicaciones} & \num{5.5} & \num{0.0} & \num{-5.5} & \num{-5.7} & \num{-100.0} \\
{49 Fiscalía General de la República} & \num{515.2} & \num{0.0} & \num{-515.2} & \num{-537.4} & \num{-100.0} \\
\bottomrule
\end{tabular}
\par\vspace{2pt}\footnotesize Fuente: elaboración propia del ITED con los analíticos del PEF aprobado 2024 y 2025. Las diferencias en millones de pesos se reportan dos veces, nominal y real; la real está en pesos de 2025 con el deflactor de 4.3\% que declara el CGPE.
\end{table}


\section{La obra pública directa se contrae}

El capítulo 6000, que es obra pública ejecutada directamente por el Gobierno
Federal y las entidades de control directo, pasa de \textbf{447,575.3 a 290,693.9
millones de pesos}: una caída real de \textbf{37.7\%}. El capítulo 5000, bienes
muebles e inmuebles, cae 43.5\%. Las inversiones financieras, en cambio, se
mantienen: 0.9\% real de caída sobre 340,623.7 millones.

\textbf{El paquete que consolida lo hace sobre la obra, no sobre la inversión
financiera.} Es una elección de composición y tiene consecuencia: la inversión
financiera sostiene balances de empresas, la obra construye activos.

\section{Quién construye}

El movimiento por ramo confirma lo que el capítulo 4 anticipó. En obra pública
directa, Marina pasa de cero a \textbf{22,058.5 millones de pesos}, mientras la
Secretaría de Infraestructura baja de 35,330.5 a 11,833.7 y Medio Ambiente de
35,343.4 a 6,682.1.

Leído junto con el traslado de la obra ferroviaria de Defensa a la Secretaría de
Infraestructura, el cuadro es de \textbf{recomposición de ejecutores más que de
niveles}: cambia quién construye antes que cuánto se construye, y en el agregado
sí baja.

\section{Lo que este capítulo no puede decir}

El paquete no publica una cartera de proyectos con montos en las piezas que este
documento tiene. El tomo de proyectos de inversión no está en la carpeta.
\textbf{No se puede, por tanto, decir qué proyectos concentran la obra de 2025}, y
la pregunta queda abierta.

\clearpage
